% 附录回归测试：固化附录编号（A/B/C）与“无空白页”行为
%
% 历史 bug（PR #12）：myappendix 在 \backmatter 之后调用时 \appendix 依赖
% \@mainmattertrue 才能生成“附录A/B/C”编号，否则附录章节无编号、且
% \chapter*{附录} 总标题会独占一页产生空白页。
%
% 本测试通过 \TYPE 显式输出各计数器的值，固化以下行为：
%   1. 附录章节编号为 A/B/C（\thechapter = \Alph{chapter}）
%   2. 附录内表格/图表/公式编号带章节前缀（表 A.1、图 B.1 等）
%   3. 页标记连续（[1]→附录A、[2]→附录B、[3]→附录C），不存在空白页
\documentclass{sduthesis}
% 测试用公共配置（不含 documentclass，由各测试文件自行声明）
\SDUSetup{
  module       = undergraduate,
  title        = {基于卷积神经网络的图像识别方法研究},
  author       = {张某某},
  studentId    = {202100100001},
  school       = {信息科学与工程学院},
  major        = {计算机科学与技术},
  supervisor   = {李某某},
  year         = {2025},
  month        = {5},
}

\input{regression-test}

\begin{document}
\START
% 正文章节（非附录），用于验证编号切换
\chapter{引言}
\TYPE{正文 chapter number: \thechapter}

% 附录必须在 \backmatter 之前调用（与 main.tex 一致）
\begin{myappendix}
\chapter{攻读学位期间发表的学术论文}

\TYPE{appendix A chapter number: \thechapter}
\TYPE{appendix A equation number: \theequation}

\begin{enumerate}
    \item 张三, 李四. 基于深度学习的图像识别算法研究[J]. 计算机科学与探索, 2024, 18(5): 123-135.
\end{enumerate}

\chapter{补充实验数据}

\section{网络各层详细参数统计}

\begin{table}[htbp]
    \centering
    \caption{网络各层参数统计}
    \begin{tabular}{cccc}
        \toprule
        \textbf{层类型} & \textbf{数量} & \textbf{参数量(M)} & \textbf{占比(\%)} \\
        \midrule
        $1\times1$ 卷积 & 12 & 0.82 & 25.6 \\
        $3\times3$ 深度卷积 & 15 & 0.34 & 10.6 \\
        \bottomrule
    \end{tabular}
    \label{tab:layer_params}
\end{table}

\TYPE{appendix B chapter number: \thechapter}
\TYPE{appendix B table number: \thetable}

\begin{equation}
    E = mc^2
\end{equation}

\chapter{代码实现补充说明}

\TYPE{appendix C chapter number: \thechapter}

\begin{figure}[htbp]
    \centering
    \rule{2cm}{1cm}
    \caption{附录示意图}
    \label{fig:appendix}
\end{figure}

\TYPE{appendix C figure number: \thefigure}

\end{myappendix}
\END
\end{document}
